% !TEX encoding = UTF-8 Unicode
\chapter{Kernel Jupyter para LisPy}
\label{chap:jupyter}
{\huge$\bigstar$}{\huge$\bigstar$} Intermedio

LisPy incluye un kernel Jupyter completo que permite usar LisPy en Jupyter Notebook,
JupyterLab y VS Code. El kernel mantiene \textbf{estado persistente} entre celdas:
variables Lisp, funciones, macros, módulos Python importados y el namespace
\texttt{\_\_main\_\_} de Python.

\section{Instalación}

\begin{lispcode}
$ cd jupyter
$ python install.py
\end{lispcode}

Esto registra el kernel \texttt{LisPy (Common Lisp + Python)} en Jupyter y configura
el resaltado de sintaxis CodeMirror para archivos \texttt{.lp}.

\begin{lispcode}
;;; Verificar instalación
jupyter kernelspec list
;;; Debe mostrar: lisPy    C:\Users\...\jupyter\kernels\lisPy
\end{lispcode}

\section{Arquitectura del Kernel}

\begin{enumerate}
\item \textbf{Proceso persistente}: El kernel lanza \texttt{lispy.exe --kernel}
      como subproceso que mantiene todo el estado entre celdas.
\item \textbf{Preprocesador híbrido}: Cada celda se preprocesa con
      \texttt{\_lp\_preprocess.py} para soportar sintaxis híbrida
      Lisp/Python (Capítulo~\ref{chap:hybrid-syntax}).
\item \textbf{Librerías auto-cargadas}: Al iniciar, el kernel carga automáticamente
      \texttt{ffi}, \texttt{continuations}, \texttt{defpyfun}, \texttt{defpyclass},
      \texttt{python-template} y \texttt{persist}.
\item \textbf{Venv automático}: El kernel detecta y usa el entorno virtual
      en \texttt{jupyter/venv/} si existe.
\end{enumerate}

\section{Uso Básico}

\begin{lispcode}
;;; Celda 1 — Definir función
(defun fib (n)
  (if (<= n 1) n (+ (fib (- n 1)) (fib (- n 2)))))
(fib 20)
\end{lispcode}

\begin{lispcode}
;;; Celda 2 — La función persiste (no hay reinicio)
(mapcar #'fib '(1 2 3 4 5 6 7 8 9 10))
\end{lispcode}

\begin{lispcode}
;;; Celda 3 — Python con estado persistente
(py:import "math")
(py:call (py:attr math "sqrt") 144)
\end{lispcode}

\begin{lispcode}
;;; Celda 4 — Símbolos auto-resueltos (sin re-importar)
math
(py:call (py:attr math "cos") 0)
\end{lispcode}

\section{Sintaxis Híbrida en Celdas}

Cada celda del notebook se procesa automáticamente con el preprocesador híbrido:

\begin{lispcode}
;;; Definir función Lisp
(defun cuadrado (x) (* x x))

;;; Usar en Python (misma celda o celda posterior)
for i in range(1, 6):
    print(f"{i}² = {cuadrado(i)}")
\end{lispcode}

\section{Librerías Pesadas (Carga Bajo Demanda)}

\begin{lispcode}
;;; Cargar NumPy (importa numpy + configura bindings)
(load "lib/numpy-fast.lp")
(numpy-linspace 0 pi 5)

;;; Cargar SymPy
(load "lib/sympy-bridge.lp")
(let ((x (sympy-symbol "x")))
  (sympy-derive (sympy-expr (sin (expt x 2))) x))
\end{lispcode}

\section{Resaltado de Sintaxis}

El kernel incluye soporte para CodeMirror 6 (JupyterLab 4+) con:
\begin{itemize}
\item Palabras clave de Common Lisp
\item Símbolos del paquete \texttt{py:}
\item Macros de lectura (\texttt{\#\_}, \texttt{\#'}, \texttt{\#(})
\item Sintaxis híbrida Python embebida
\end{itemize}

También está disponible la extensión VS Code en \texttt{lp-language/} con soporte
completo de TextMate grammar para archivos \texttt{.lp}.

\section{Notebook de Demostración}

\begin{lispcode}
cd jupyter
jupyter notebook LisPy-Demo.ipynb
\end{lispcode}

El notebook \texttt{LisPy-Demo.ipynb} contiene 16 secciones que cubren todas las
capacidades de LisPy, desde aritmética básica hasta machine learning con
scikit-learn.

\section{Proyectos de Ejemplo}

La carpeta \texttt{jupyter/projects/} contiene 6 proyectos completos autocontenidos:
\begin{itemize}
\item \textbf{01 — Análisis de Datos}: Pipeline completo con pandas, matplotlib, GROUP BY
\item \textbf{02 — Machine Learning}: Clasificación, regresión, pipelines scikit-learn
\item \textbf{03 — Matemáticas Híbridas}: \texttt{defmath}, SymPy, gradiente, optimización
\item \textbf{04 — Trading Engine}: CLOS, backtesting, Sharpe ratio
\item \textbf{05 — Procesamiento de Señales}: FFT, filtros, streaming con generadores
\item \textbf{06 — Robot 2-DOF}: Lagrange simbólico, solve\_ivp, control PD
\end{itemize}

\section{Extensión VS Code}

\begin{lispcode}
;;; Instalar extensión VS Code para archivos .lp
cd lp-language
code --install-extension lp-language-0.0.1.vsix
\end{lispcode}

La extensión proporciona:
\begin{itemize}
\item Resaltado de sintaxis (TextMate grammar + Python injection)
\item Comentarios de línea (\texttt{;}) y bloque (\texttt{\#| |\#})
\item Bracket matching, auto-closing, folding
\item Soporte CodeMirror para notebooks integrados en VS Code
\end{itemize}
